% =============================================================================
% SYSTEM INSTRUCTION TO LLM:
% This file is a PROMPT TEMPLATE.
% You MUST generate a FULL, FINAL, COMPILABLE LaTeX presentation
% based on this template.
%
% CRITICAL RULES:
% 1. DO NOT change anything in the Preamble.
% 2. You MUST delete all [STYLE REF] frames and the AI INSTRUCTION frame.
% 3. You MUST replace all placeholder texts with real content.
% 4. You MUST preserve the collegeBeamer theme and visual style: this deck is
%    SWISS / INTERNATIONAL TYPOGRAPHIC — bold sans (TeX Gyre Heros), stark
%    near-black on white, flush-left, square geometry. NEVER introduce colors:
%    everything is maincolor (near-black) plus gray tints.
% 5a. Keep each slide's body text to at most 8--10 lines.
% 5b. Total deck length: 12--16 slides INCLUDING the title, outline, section
%     divider and Q&A pages; use 3--4 sections for short source documents.
% 6. NEVER use \includegraphics. No external image files exist. Draw every chart,
%    diagram and table with tikz / pgfplots / booktabs instead.
% 7. Only these packages are available (already loaded in the preamble): collegebeamer,
%    tikz, pgfplots, amsmath, amssymb, pifont, booktabs, colortbl, multirow.
%    Do NOT add \usepackage lines.
% 8. NO pie charts (no pie package is available) -- for composition/share data use
%    a stacked ybar, a horizontal xbar, or a booktabs table instead.
% 9. Tables: max 5 columns; wrap long text columns in p{..\textwidth}; use
%    \footnotesize; never let a table exceed \textwidth.
% 10. Write ALL slide text in English by default, even for non-English source
%     documents; switch languages only when the caller's requirements ask for
%     it, and then set \date in that language too (\today is an English date).
% 11. Charts are MONOCHROME: bars filled solid maincolor, secondary series in
%     black!45 dashed; no axis box (axis line style={draw=none}), horizontal
%     gridlines only. Copy the axis options from the examples verbatim.
% =============================================================================

\documentclass[aspectratio=169,10pt]{beamer}

% -----------------------------------------------------------------------------
% 1. Preamble & Theme (DO NOT CHANGE)
% -----------------------------------------------------------------------------
\usepackage[zh,bjtu]{collegebeamer}

\usepackage{tikz}
\usetikzlibrary{positioning}
\usepackage{pgfplots}
\pgfplotsset{compat=1.17}
\usepackage{amsmath,amssymb}
\usepackage{pifont}
\usepackage{booktabs}
\usepackage{colortbl}
\usepackage{multirow}

% --- Font Setup ---
% Latin -> TeX Gyre Heros (Helvetica clone), CJK -> FandolHei; both set by the
% theme. The deck is sans-first: do not switch families.
% Disable collegeBeamer auto TOC insertion (section dividers are hand-built)
\AtBeginSection[]{}

% -----------------------------------------------------------------------------
% 1.1 Title page (DO NOT CHANGE)
% -----------------------------------------------------------------------------
% The theme provides the Swiss title page: thick 12x4mm near-black bar, huge
% bold flush-left title anchored top-left, tight metadata stack, small date
% pinned bottom-left. Subtitle/author lines vanish when their fields are empty,
% so the page stays balanced even with the default empty metadata.
% Frames placed before the first \section print no section eyebrow (no stray "0").

% -----------------------------------------------------------------------------
% 2. Metadata (LLM MUST FILL)
% -----------------------------------------------------------------------------
\title[Short Title]{Full Presentation Title}
% Caller fills these three fields. Leave empty when unknown; never invent placeholders.
\subtitle{}
\author{}
\institute{}
\date{\today}

\begin{document}

% =============================================================================
% FRAME 1: TITLE PAGE
% =============================================================================
\maketitle

% =============================================================================
% FRAME 2: MANUAL TABLE OF CONTENTS
% Swiss list: zero-padded bold numerals, generous leading, flush left.
% =============================================================================
\begin{frame}{Outline}
\vspace{1mm}
\large
\renewcommand{\baselinestretch}{1.8}\selectfont

\textbf{01}\hspace{1.2em}Title of Part One

\textbf{02}\hspace{1.2em}Title of Part Two

\textbf{03}\hspace{1.2em}Title of Part Three

\textbf{04}\hspace{1.2em}Title of Part Four

\textbf{05}\hspace{1.2em}Summary \& Q\&A
\end{frame}

% =============================================================================
% SECTION 1 (EXAMPLE)
% IMPORTANT: define \section{...} before any content frames of that section
% so the eyebrow line under each frametitle is complete (no "0").
% =============================================================================
\section{Title of Part One}

% =============================================================================
% FRAME 3: SECTION DIVIDER (EXAMPLE: SECTION 1)
% LLM MUST generate one such page right after EACH \section{...}.
% White page, giant ghost numeral in 8% black, thick near-black bar, bold
% section name. Update the numeral (01 -> 02, 03 ...) to the section number.
% NO full-bleed color pages anywhere in this deck.
% =============================================================================
\begin{frame}[plain,t]
\vspace*{0.18\paperheight}
{\fontsize{100}{100}\selectfont\bfseries\textcolor{black!8}{01}\par}
\vspace{-12mm}
\textcolor{maincolor}{\rule{12mm}{4mm}}\par
\vspace{8pt}
{\fontsize{20}{24}\selectfont\bfseries\textcolor{maincolor}{Title of Part One}\par}
\end{frame}

% =============================================================================
% CONTENT FRAME TEMPLATE (LLM MUST USE THIS STRUCTURE)
% =============================================================================
% Content page example 1: Text-left, figure-right layout
\begin{frame}{Content Example: Text Left, Figure Right}
\framesubtitle{} % keep empty; the theme fills the section eyebrow automatically
\small
\begin{columns}[T]
\column{0.55\textwidth}
Write the core analysis content on the left side here.
This area is used for background description, problem definition,
logical reasoning, or methodology explanation.
Text should be concise and conclusion-oriented; avoid vague descriptions.

\vspace{0.6em}
\textbf{Key Takeaway:}
State the central message of this slide in one sentence.

\column{0.40\textwidth}
\centering
% Side panel: FEW categories only (4--6 short labels). Compact width. No xlabel.
% Swiss chart: solid near-black bars, no axis box, horizontal gridlines only.
% For 12-month / many-label series use the full-width chart on the next example page.
\begin{tikzpicture}
\begin{axis}[
    width=5.4cm,
    height=3.6cm,
    ybar,
    bar width=10pt,
    enlarge x limits=0.25,
    ymin=0, ymax=100,
    symbolic x coords={Q1,Q2,Q3,Q4},
    xtick=data,
    ylabel={Sample index},
    tick label style={font=\tiny},
    label style={font=\scriptsize},
    axis line style={draw=none},
    xtick style={draw=none},
    ytick style={draw=none},
    ymajorgrids=true,
    grid style={black!14},
]
\addplot[fill=maincolor, draw=none] coordinates {
    (Q1,42) (Q2,55) (Q3,61) (Q4,78)
};
\end{axis}
\end{tikzpicture}
\end{columns}
\end{frame}

% Content page example 2: Text-top, figure-bottom layout
\begin{frame}{Content Example: Text Top, Figure Bottom}
\framesubtitle{} % automatically inherits the section eyebrow
\small

% Upper text area
When a series has many categories (for example 12 months), use this
full-width layout below -- not the narrow side-panel chart.
Keep the text concise so the wide chart remains readable.

\vspace{0.2em}

% Lower figure: MANY categories (12 months) → MUST be wide + rotated x labels.
% Monochrome line chart: primary series solid maincolor, secondary black!45 dashed.
\centering
\begin{tikzpicture}
\begin{axis}[
    width=11.5cm, height=3.7cm,
    enlarge x limits=0.04,
    x tick label style={rotate=45, anchor=east, font=\tiny},
    ylabel style={font=\scriptsize},
    xlabel style={font=\scriptsize, yshift=-2pt},
    ymin=1900, ymax=2900,
    symbolic x coords={Jan,Feb,Mar,Apr,May,Jun,Jul,Aug,Sep,Oct,Nov,Dec},
    xtick=data,
    ylabel={Sample level},
    xlabel={Month (illustrative)},
    legend style={at={(0.02,0.98)}, anchor=north west, font=\tiny, draw=none, fill=none},
    yticklabel style={font=\tiny},
    axis line style={draw=none},
    xtick style={draw=none},
    ytick style={draw=none},
    ymajorgrids=true,
    grid style={black!14},
]
\addplot[
    very thick,
    maincolor,
    mark=square*,
    mark options={fill=maincolor, scale=0.6},
] coordinates {
    (Jan,2040) (Feb,2085) (Mar,2160) (Apr,2330) (May,2345) (Jun,2320)
    (Jul,2420) (Aug,2505) (Sep,2630) (Oct,2740) (Nov,2655) (Dec,2620)
};
\addplot[
    thick,
    dashed,
    black!45,
    mark=square*,
    mark options={fill=white, draw=black!45, scale=0.55, solid},
] coordinates {
    (Jan,2010) (Feb,2050) (Mar,2120) (Apr,2280) (May,2300) (Jun,2290)
    (Jul,2380) (Aug,2460) (Sep,2580) (Oct,2700) (Nov,2620) (Dec,2590)
};
\legend{Series A (sample), Series B (sample)}
\end{axis}
\end{tikzpicture}

\vspace{0.1em}
\centerstate{
    \tiny \emph{Source: data compiled from the source document.}
}
\end{frame}


% Content page example 3: Figure-left, text-right layout
\begin{frame}{Content Example: Figure Left, Text Right}
\framesubtitle{}
\small
\begin{columns}[c] % [c] = vertically centered; use [T] for top alignment if text is long

% Left: real TikZ structure diagram (layered pipeline) — SQUARE corners only,
% near-black strokes, gray-tint fills. Never use rounded corners in this deck.
\column{0.45\textwidth}
    \centering
    \begin{tikzpicture}[
        font=\scriptsize,
        box/.style={
            draw=maincolor, line width=0.8pt, fill=white,
            minimum height=0.62cm, text width=1.35cm, align=center
        },
        layer title/.style={font=\tiny\bfseries, text=maincolor}
    ]
        % Data layer
        \fill[black!6] (-0.25,1.05) rectangle (4.85,2.45);
        \node[layer title, anchor=north west] at (-0.15,2.40) {Data layer};
        \node[box] (s1) at (0.7,1.55) {Sensors};
        \node[box] (s2) at (2.5,1.55) {Ingest};
        \node[box] (s3) at (4.2,1.55) {Store};
        % Decision layer
        \fill[black!6] (-0.25,-0.95) rectangle (4.85,0.70);
        \node[layer title, anchor=north west] at (-0.15,0.65) {Decision layer};
        \node[box] (p1) at (0.7,-0.15) {Feature};
        \node[box] (p2) at (2.5,-0.15) {Model};
        \node[box] (p3) at (4.2,-0.15) {Act};
        % Edges
        \draw[->, thick, maincolor] (s1) -- (s2);
        \draw[->, thick, maincolor] (s2) -- (s3);
        \draw[->, thick, maincolor] (s2) -- (p1);
        \draw[->, thick, maincolor] (p1) -- (p2);
        \draw[->, thick, maincolor] (p2) -- (p3);
        \draw[->, thick, maincolor, dashed]
            (p3.south) -- ++(0,-0.28) -| (p1.south)
            node[pos=0.25, below, font=\tiny, text=maincolor] {feedback};
    \end{tikzpicture}

% Right: text analysis area (50\% width)
\column{0.50\textwidth}
    \textbf{Figure / Object Analysis:}

    Interpret the visual element on the left. Suitable for presenting:
    \begin{itemize}
        \item Structure of experimental apparatus
        \item Module details of an algorithm or model
    \end{itemize}

    \vspace{0.35em}
    \textbf{Key Findings:}
    \begin{enumerate}
        \item \ctextbf{Observation 1}: the left figure shows\ldots
        \item \ctextbf{Observation 2}: results differ from expectations\ldots
    \end{enumerate}
\end{columns}
\end{frame}

% =============================================================================
% [STYLE REF] TABLE PAGE (DELETE IN FINAL OUTPUT)
% =============================================================================
\begin{frame}{[STYLE REF] Table Style Example}
\framesubtitle{}
\small
\centering
% Swiss table: bold sans header, generous row height (arraystretch 1.5),
% horizontal booktabs rules ONLY — never add vertical rules (no | in the spec).
% Text columns left-aligned, number columns right-aligned.
% No \rowcolors zebra striping on short tables (fewer than ~5 data rows).
\renewcommand{\arraystretch}{1.5}
\footnotesize
\begin{tabular}{l r r l}
\toprule
\textbf{Metric} & \textbf{A} & \textbf{B} & \textbf{Conclusion} \\
\midrule
Item 1 & 12 & 15 & Room for improvement \\
Item 2 & 20 & 18 & Stable \\
\bottomrule
\end{tabular}
\end{frame}

% =============================================================================
% [STYLE REF] TIKZ FLOW PAGE (DELETE IN FINAL OUTPUT)
% =============================================================================
\begin{frame}{[STYLE REF] Flowchart Style Example}
\framesubtitle{}
\centering
% Swiss flow: sharp rectangles, near-black strokes, white or gray-tint fills.
\begin{tikzpicture}[
node distance=0.9cm and 0.7cm,
process/.style={rectangle,draw=maincolor,line width=0.8pt,fill=black!6,
minimum height=0.95cm,text width=1.9cm,align=center}
]
\node[process] (a) {Collect};
\node[process,right=of a] (b) {Analyze};
\node[process,right=of b] (c) {Decide};
\node[process,below=of b] (d) {Revise};
\draw[->,thick,maincolor] (a)--(b);
\draw[->,thick,maincolor] (b)--(c);
\draw[->,thick,maincolor] (b)--(d)
    node[midway,right,font=\tiny,text=maincolor] {gap};
\draw[->,thick,maincolor,dashed] (d.west) -| (a.south)
    node[near end,left,font=\tiny,text=maincolor] {retry};
\end{tikzpicture}
\end{frame}

% =============================================================================
% [STYLE REF] TWO-COLUMN ANALYSIS (DELETE IN FINAL OUTPUT)
% =============================================================================
\begin{frame}{[STYLE REF] Two-Column Analysis Example}
\framesubtitle{}
\small
\begin{columns}
\column{0.52\textwidth}
Left column for argumentation, definitions, and derivations.

\column{0.43\textwidth}
Right column for action items, illustrations, or comparisons.
\end{columns}
\end{frame}

% =============================================================================
% AI INSTRUCTION PAGE (DELETE IN FINAL OUTPUT)
% =============================================================================
\begin{frame}[allowframebreaks]{AI GENERATION INSTRUCTIONS}
\framesubtitle{}
You MUST:
\begin{itemize}
\item Delete all [STYLE REF] frames and this instruction frame
\item Replace all placeholders with real content
\item Insert a section divider page (giant ghost numeral + thick bar + bold
      section name, as in Frame 3) right after each \texttt{\textbackslash section\{...\}}
\item Define \texttt{\textbackslash section\{...\}} before the first content frame of that section
\item Use the columns layout for text+side-figure pages; use the full-width
      pattern (Example 2) whenever a chart has more than 6 categories
\item Keep everything monochrome: maincolor (near-black) and gray tints only;
      solid maincolor bars, black!45 dashed secondary lines, no other colors
\item Keep all geometry square: no rounded corners, no circles, no color fills
      other than white / gray tints
\item Keep information density high (no empty slides)
\item Output FULL LaTeX code only
\end{itemize}
\end{frame}

% =============================================================================
% FINAL PAGE
% =============================================================================
\QApage

\end{document}
